% Actual full derivations for the selected Turner 2008 chapter.
\section{Nonmonotonic causal theories: proofs}
\label{sec:annex-p16-nonmonotonic-causal-theories}
Use the conventions of \cref{sec:p16-nonmonotonic-causal-theories}.
These are independent derivations at the stated scope, not claimed readings
of every proof in the source's bibliography
\cite[pp.~759--776]{Turner2008NonmonotonicCausalLogic}.

\begin{annexitem}{annex:p16.nct-basic}{thm:p16-nct-basic}{Truth theories and the nonmonotonic example}
\begin{proof}
A unique reduct model makes consequence exactly its truth theory.
Conversely, the stated equivalence for all formulas includes \(\bot\):
since it is false in \(I\), \(T^I\) does not entail it and has a model.
Every true atom \(c=I(c)\) is entailed; any reduct model therefore agrees
with \(I\) at every coordinate and equals it. The empty signature is included.
An activated rule's head belongs to \(T^I\), hence holds in its model.
This proves each material implication. Its converse fails for the empty
theory on one Boolean constant: both worlds satisfy all zero implications,
but neither is the unique model of its empty reduct.

For the three-value example the reduct at value 1 is \(\{p=1\}\);
at either other value it is empty with three models. Adding the self-rule
for value 2 changes its reduct to \(\{p=2\}\) while value 3's remains empty.
This proves all asserted model counts.

The isolated head-selector in the printed p.~763 definition would alter
the semantics. Take \(p\Leftarrow q,q\Leftarrow q,\neg p\Leftarrow\neg p\).
At \(I=\{\neg p,q\}\), true-head selection gives \(\{\neg p,q\}\), uniquely
satisfied by \(I\). Correct body activation gives \(\{p,q,\neg p\}\),
which has no model. Thus the inconsistent display cannot be substituted
for the source's p.~759 definition and its later calculations.
\end{proof}
\end{annexitem}

\begin{annexitem}{annex:p16.nct-extensions}{thm:p16-nct-extensions}{Filtering and conservative extensions}
\begin{proof}
A true constraint body activates \(\bot\); a false one activates nothing.
The reduct is respectively unsatisfiable or unchanged, proving exact filtering.

Old reduct heads contain no fresh \(d\). With \(d\equiv F\Leftarrow\top\),
each old reduct model \(J\) extends exactly once, with \(d\) the truth value
of \(F\) at \(J\). This is a bijection of reduct models.
The actual candidate must be that prescribed extension, and uniqueness
is equivalent to uniqueness of the old reduct.

For the value-default construction, fix an old \(I\).
Incompatibility allows at most one true \(F_i\), \(i\geq2\).
If \(F_j\) holds, the reduct forces \(d=v_j\); any different actual value
fails that head, and actual \(v_1\) additionally activates its conflicting
self-head. At the prescribed \(v_j\), the new heads simply fix \(v_j\).
If none holds, actual \(v_1\) activates the only self-head and fixes it.
A different actual value has no activated head for \(d\); if the old
reduct is satisfiable and \(r>1\), varying \(d\) gives more than one model,
so that candidate is not explanatory. An unsatisfiable old reduct admits
no explanatory extension in any case. For \(r=1\), the sole value is already
fixed by its domain, and its self-rule is harmless.
In every prescribed case the new reduct is the old reduct plus a fixed fresh
value, proving the bijection and default claim.
If two \(F_i\) hold after removing incompatibility, their different heads
conflict, proving why the hypothesis matters.
\end{proof}
\end{annexitem}

\begin{annexitem}{annex:p16.nct-strong}{thm:p16-nct-strong}{SE characterization and explicit contexts}
\begin{proof}
A union's SE-model set is the intersection, since its reduct is the union
of activated heads. An \(I\) is explanatory iff \((I,I)\) is an SE-model
and no \((I,J)\), \(J\ne I\), is one. Equality of SE-model sets therefore
preserves models after intersection with every context, proving sufficiency.

For necessity suppose \((I,J)\) belongs to the SE-models of \(S\) but not
\(T\), exchanging names for the opposite difference.
If \(I\not\models T^I\), add
\(R=\{c=I(c)\Leftarrow\top:c\in\sigma\}\).
Its heads fix exactly \(I\). This makes \(I\) explanatory for \(S\cup R\),
but it still violates \(T^I\), so not for \(T\cup R\).
Otherwise \(I\models T^I\) and \(J\not\models T^I\); thus \(J\ne I\).
Choose \(c_*\) on which they differ. Let \(E\) contain
\(c_*=I(c_*)\lor c_*=J(c_*)\), and for every constant \(c\) both
\[
 (c_*=I(c_*))\supset(c=I(c)),\qquad
 (c_*=J(c_*))\supset(c=J(c)).
\]
Its models are exactly \(I,J\): the disjunction chooses a branch and the
remaining formulas fix all its coordinates. They remain finite formulas
on an infinite signature. Let \(R=\{F\Leftarrow\top:F\in E\}\).
At \(I\), \(S\cup R\)'s reduct admits both \(I,J\), while \(T\cup R\)'s
admits exactly \(I\). Hence the contexts separate them. These are set-sized
contexts on the original signature; no finite-signature assumption entered.
If no interpretations exist, all SE-model/model sets are empty and the
assertion is vacuous.
Equivalent bodies activate the same heads at every \(I\); equivalent heads
hold at the same \(I,J\). Rulewise equivalent replacement thus preserves
SE-models and strong equivalence.

For the source example, \(S=\{p\Leftarrow q,q\Leftarrow p\}\) activates
both heads at \(p,q\). At any other world it either activates none or a
violated head. \(T=\{p\Leftarrow p,q\Leftarrow q\}\) uniquely fixes both
values only at \(p,q\); missing positives have no negative explanation.
Adding the negative \(p\) self-rule to \(T\) additionally permits
\(\neg p,q\), whose reduct fixes those two literals.
At that world \(S\)'s first rule also activates \(p\), producing a
contradiction. The worlds with \(q=f\) still lack a negative \(q\) head.
These checks exhaust all four interpretations.
\end{proof}
\end{annexitem}

\begin{annexitem}{annex:p16.nct-completion}{thm:p16-nct-completion}{Every nontrivial coordinate must be supported}
\begin{proof}
In an explanatory \(I\), no constraint activates.
Every activated atom head is true. Conversely, suppose the true nontrivial
atom \(A=(c=I(c))\) is not an activated head.
Any activated atom about \(c\) must equal \(A\), because it holds in \(I\).
There are therefore no heads about \(c\). Changing just \(I(c)\) to
another domain value preserves every reduct head and gives another model,
contradicting uniqueness. Some body for \(A\) must hold.
For every nontrivial value atom, truth is therefore equivalent to activation
of one of its finitely many bodies. This is completion.

Conversely, completion ensures false constraint bodies and true activated
atom heads, including trivial ones forced by domains.
For each nontrivial constant the equation for its actual value activates
a head fixing that value. Any reduct model agrees at each such constant;
singleton domains fix the rest. Hence \(I\) is its unique model.
The proof applies to set-sized signatures; finiteness of each head-list
makes the formulas legitimate, and no graph acyclicity was used.
On a singleton domain the empty theory has one world; the forbidden
extra equation \(c=0\equiv\bot\) excludes it.
For Boolean \(p\Leftarrow p\), the negative-value equation is
\(\neg p\equiv\bot\), forcing true. The positive program's least model is
empty, while its Clark biconditional is tautological, proving the comparison.
\end{proof}
\end{annexitem}

\begin{annexitem}{annex:p16.nct-complexity}{thm:p16-nct-complexity}{Certificates and polynomial reductions}
\begin{proof}
An explicitly encoded interpretation has polynomial length, and evaluating
finite formula parse trees is polynomial. Well-formedness, domains and
value membership are polynomial checks.

Definite completion has polynomial size: each body's occurrence belongs
to one head-list, every listed nontrivial value has one equation, and
domain constraints use at most quadratically many exclusions per domain.
Guessing a satisfying value assignment proves NP membership.
For SAT hardness, given Boolean \(F\), take all value self-rules
\(c=t\Leftarrow c=t,c=f\Leftarrow c=f\), and
\(\bot\Leftarrow\neg F\). The self-rules fix every world's reduct;
the constraint retains exactly the satisfying assignments of \(F\).
The theory is definite and polynomial. The existing SAT benchmark gives
NP completeness.

Non-strong-equivalence has a polynomial certificate \(I,J\): evaluate
each body at \(I\), and each activated head at \(I,J\), then identify
different SE membership. This proves coNP membership of strong equivalence.
For hardness compare \(\{\bot\Leftarrow F\}\) to the empty theory.
The former has an SE-pair iff \(I\not\models F\), the latter every pair.
Their SE-models coincide iff \(F\) is unsatisfiable, reducing UNSAT.

For ordinary non-equivalence, guess an orientation and \(I\) explanatory
for \(T\) but not \(S\). The latter is witnessed either by \(I\)'s violating
a head of \(S^I\), or by another \(J\ne I\) satisfying that reduct.
For \(T\), require \(I\models T^I\) and universally check
\(K\models T^I\Rightarrow K=I\).
This is an existential block of polynomial certificates followed by one
universal interpretation block and a polynomial predicate, so
non-equivalence is in \(\Sigma_2^p\) and equivalence in \(\Pi_2^p\).

For hardness take \(\forall X\exists Y\,F(X,Y)\) and fresh Boolean \(z\).
Set \(B_Y=\bigwedge_{y\in Y}\neg y\), empty conjunction \(\top\), and
\(H=(z\land B_Y)\lor(\neg z\land F(X,Y))\).
Let \(T_F\) contain all exogeneity rules for \(X\),
\(\bot\Leftarrow\neg z\), and \(H\Leftarrow z\).
Compare with \(U=\{\bot\Leftarrow\top\}\), which has no models.
An explanatory candidate of \(T_F\) must have \(z=t,Y=f\).
At any such fixed \(X\)-assignment its reduct has that candidate and,
as its only alternatives, worlds with the same \(X\), \(z=f\), and
\(F(X,Y)\) true. Thus the candidate is unique iff no \(Y\)-witness
exists at that \(X\). Consequently \(T_F,U\) are equivalent exactly
when \(\forall X\exists Y\,F\) is true.
The construction is polynomial and works also for empty variable blocks.
The existing fixed-alternation QBF benchmark gives \(\Pi_2^p\) hardness,
completing all three classifications under the stated encoding.
\end{proof}
\end{annexitem}

\begin{annexitem}{annex:p16.nct-action-models}{thm:p16-nct-action-models}{Exhaustive action calculations}
\begin{proof}
Each theory is definite when negative Boolean heads are false-value atoms.
Apply the proved completion theorem.

For the bowl abbreviate truth values by \(L_0,R_0,P_0,L_1,R_1,P_1,A,B\).
Initial/action self-completions are tautological; the initial static rule
additionally requires \(L_0\ne R_0\Rightarrow P_0\).
The successor handle equations are
\(L_1\equiv(L_1\land L_0)\lor A\) and
\(\neg L_1\equiv\neg L_1\land\neg L_0\), and symmetrically for \(R\).
If \(A\) or \(L_0\) is true, respectively the positive or negative equation
forces \(L_1=t\). If both are false, the positive equation forces false.
Thus \(L_1=L_0\lor A\), and \(R_1=R_0\lor B\).
The spill equations are
\(P_1\equiv(P_1\land P_0)\lor(L_1\ne R_1)\) and
\(\neg P_1\equiv\neg P_1\land\neg P_0\), giving
\(P_1=P_0\lor(L_1\ne R_1)\) by the same cases.
Substitute the all-false initial values and each of the four action
assignments. This yields exactly the four listed successors and excludes
the unsupported right-up alternative after only a left lift.

For the cycle, fix the initial value and action.
If absent, the only support for any successor is inertia, so it must equal
the initial value. If present, the indicated effect forces the next cyclic
value; its equation has an active effect body, and every other value
equation is false on both sides. Initial/action exogeneity permits each
of three initial values and two action values, giving exactly six.

For the coin, completion at each successor value \(v\) is
\[
 coin_1=v\equiv(coin_1=v\land coin_0=v)\lor(coin_1=v\land toss_0).
\]
Untossed, only the initial value has support. Tossed, either successor
satisfies its equation, and the other value equation is false on both
sides. This gives two untossed and four tossed worlds, without weights.

For location, every occurring \(put(v)\) forces its value.
Distinct concurrent puts give incompatible reduct heads, hence no model.
With none, inertia supports exactly the old location. With one, its head
fixes the new location; inertia either agrees or is inactive.
All initial locations and these \(k+1\) action patterns are supported
by their exogeneity heads, giving \(k(k+1)\) models.

For dominos, no down domino can rise: there is no positive support except
inertia from its being up previously. Initially all are up.
Without a tip, induction on time excludes a first fall: it would need
either the absent direct effect or a preceding domino's earlier fall.
With a tip the first domino falls at time 1.
A fall of domino \(d\) between \(d-1,d\) triggers domino \(d+1\) at
time \(d+1\); before then it cannot fall because its only noninertial cause
has not activated. Once fallen it cannot rise.
Induction proves \(up(d)_t\) iff \(d>t\).
Each asserted value has an initial, direct, chain or persistence cause,
and no contradictory head activates. Both histories satisfy completion,
and the induction exhausts all possibilities, including the empty chain
family for \(k=1\).

For the pendulum, fix a current value and hold value.
A hold activates the head preserving the current value; choosing the
opposite would also activate a conflicting change-head.
Without a hold, keeping the value has no supporting head: the corresponding
change-body requires the opposite old value. Flipping activates the
appropriate change-head. These facts force exactly the asserted successor.
Induct through the horizon. All two initial choices and \(2^n\) hold
sequences are exogenously supported and each has one history,
giving \(2^{n+1}\), including horizon zero.
\end{proof}
\end{annexitem}

\begin{annexitem}{annex:p16.nct-cplus-paths}{thm:p16-nct-cplus-paths}{Local states and complete paths}
\begin{proof}
At horizon one the displayed translation is exactly the two source causal
theories. At zero, exogeneity fixes every fluent coordinate and no action
exists. The cycle admits all three values. The bowl admits exactly states
satisfying its static material condition, since a violating state activates
both the spill head and its false exogeneity head.

At any candidate history, all activated heads are value atoms at one time.
By completion, explanation requires true heads and support for every
nontrivial actual coordinate.
Initial and action values are supported by exogeneity. Successor support
is exactly that of the local one-step theory: same-time static rules,
preceding-time conditional effects and the declared inertia.
No rule in this fragment heads an earlier fluent using a later condition.

Hence an explanatory global history gives an admissible initial state and
an explanatory adjacent state/action/state triple at every step.
Conversely concatenate a finite path of those triples. Its initial/actions
have the prescribed support and each successor has precisely its local
support; static conditions hold at every time. Global completion holds,
so the history is explanatory. This proves both directions for every
horizon; zero was separately checked. The single-step outcomes already
proved give the asserted edge descriptions.
The argument is for this explicit local definite fragment, not an invented
translation of the entire C+ language.
\end{proof}
\end{annexitem}

\begin{annexitem}{annex:p16.nct-default}{thm:p16-nct-default}{The grounded default fixed point}
\begin{proof}
For \(E=\operatorname{Th}(I)\), a justification \(\psi\) is consistent
with \(E\) exactly when \(I\models\psi\): otherwise completeness includes
its negation, while if true \(I\) is a common model.
Empty prerequisites mean the grounded default operator adds exactly
the heads of \(T^I\), then classical closure.
Thus \(\Gamma_E=\operatorname{Th}(T^I)\), and \(E=\Gamma_E\) is exactly
the truth/consequence equivalence established in the first proof.

For the example, a consistent extension must choose \(q\) or \(\neg q\):
if neither is present, both justifications are consistent and generate
opposed conclusions. An inconsistent candidate rejects every justification
and generates only tautologies, so is not an extension.
Choosing \(q\) activates \(p,q\), yielding exactly their complete theory.
Choosing \(\neg q\) activates only that head, yielding exactly its theory.
Both satisfy their activation and closure equations.
The second leaves \(p\) unsettled and has two worlds as models, so is
not the truth theory of one interpretation. The first fixes both values,
proving the full extension/model comparison.
\end{proof}
\end{annexitem}

\begin{annexitem}{annex:p16.nct-higher-order}{thm:p16-nct-higher-order}{Typed uniqueness and Boolean elimination}
\begin{proof}
Fix the object domains and nonexplainable symbols.
For any alternative \(m\), each instance in \(T^*(m)\) tests the body
at the actual assignment and the head at that alternative.
Universal closure covers every free-variable assignment. Thus it
precisely expresses satisfaction of all activated heads.

Substituting \(m=N\) in \(\forall m(T^*(m)\equiv m=N)\) requires
the actual assignment to satisfy them. At every different \(m\) the
equivalence requires failure of that satisfaction. Conversely uniqueness
gives these two requirements and hence the equivalence at every alternative.
Componentwise extensional equality compares interpreted objects, functions
and predicates of the declared sorts.

In the Boolean overlap, alternatives are exactly interpretations of the
same propositional signature, and satisfaction of \(T^*(m)\) is
\(m\models T^I\). The two unique-model conditions coincide.
There are \(2^{|N|}\) Boolean assignments. By the truth definition of
universal quantification, the formula equals the conjunction of its
instances at all assignments. Each instance is ordinary propositional.
The construction gives at most an exponential multiple of input size,
not a lower bound or an optimality claim about every possible translation.
\end{proof}
\end{annexitem}

\begin{annexitem}{annex:p16.nct-ucl}{thm:p16-nct-ucl}{Universal-causation embedding}
\begin{proof}
In the translated theory, bodies and heads are nonmodal.
Bodies are evaluated at the distinguished \(I\); each active head must
hold at every \(J\in\mathcal S\). Consequently
\[
 (I,\mathcal S)\models\{\psi\supset\mathsf C\phi:\phi\Leftarrow\psi\in T\}
 \quad\Longleftrightarrow\quad I\in\mathcal S\subseteq\operatorname{Mod}(T^I).
\]
If the reduct models are exactly \(\{I\}\), the only such structure is
\((I,\{I\})\), and it satisfies the translated theory.
Conversely uniqueness of that \(I\)-model first gives \(I\models T^I\).
Any other reduct model \(J\ne I\) would give another \(I\)-model
\((I,\{I,J\})\), contradicting uniqueness.
The proof works for arbitrary set-sized signatures, but uses nonmodal
translated heads/bodies; arbitrary nested UCL formulas cannot be treated
as that reduct without further definitions.
\end{proof}
\end{annexitem}
